\setcounter{section}{3}


  \zwischenueberschrift{Übungsaufgaben}


\inputaufgabe


{}


{


Zeige, dass die nicht-leeren
\definitionsverweis {Zariski-offenen}{}{}
Teilmengen auf der
\definitionsverweis {affinen Geraden}{}{}
${\mathbb A}^{1}_{ K }$ genau die
\zusatzklammer {maximalen} {} {}
Definitionsbereiche von
\definitionsverweis {rationalen Funktionen}{}{}
sind.


}


{} {}


\inputaufgabe


{}


{


Es sei $K$ ein unendlicher
\definitionsverweis {Körper}{}{}
und es sei


\mavergleichskette


{\vergleichskette


{ P
}


{ \in }{ K[X]
}


{  }{
}


{    }{
}


{   }{
}


}
{}{}{}
ein nichtkonstantes Polynom. Zeige, dass die durch $P$ definierte Funktion
\maabb {P} { K } { K
} {}
unendlich viele Werte annimmt.


}


{} {}


\inputaufgabe


{}


{


\aufzaehlungzweiabc{Skizziere die reellen Nullstellengebilde zu


\mathl{Y^n-X^n}{.}
}{Bestimme die Verschwindungsideale zu den affin-algebraischen Mengen


\mavergleichskette


{\vergleichskette


{ V_n
}


{ \subseteq }{ {\mathbb A}^{2}_{ \R }
}


{  }{
}


{    }{
}


{   }{
}


}
{}{}{,}
die aus allen Geraden durch den Nullpunkt und durch die Eckpunkte eines regulären $n$-Ecks
\zusatzklammer {mit $(1,0)$ als einem Eck} {} {}
bestehen.
}


}


{} {}


\inputaufgabe


{}


{


Man beschreibe eine Abbildung
\maabbdisp {\varphi} { {\mathbb A}^{1}_{ K } } { {\mathbb A}^{1}_{ K }
} {,}
die bezüglich der
\definitionsverweis {Zariski-Topologie}{}{} stetig ist, die aber nicht durch ein Polynom gegeben ist.


}


{} {}


\inputaufgabe


{}


{


Es sei


\mavergleichskettedisp


{\vergleichskette


{V
}


{ \subseteq} { { {\mathbb A}_{  {\mathbb C} }^{  n   } }
}


{ =} { {\mathbb C}^n
}


{ } {
}


{ } {
}


}
{}{}{}
eine
\definitionsverweis {affin-algebraische}{}{}
Menge. Zeige, dass unter der Identifizierung


\mavergleichskettedisp


{\vergleichskette


{ {\mathbb C}^n
}


{ =} { \R^{2n}
}


{ } {
}


{ } {
}


{ } {
}


}
{}{}{}
die Teilmenge $V$ auch eine affin-algebraische Menge des ${ {\mathbb A}_{ \R }^{  2n  } }$ ist. Zeige, dass die Umkehrung nicht gilt.


}


{} {}


\inputaufgabe


{}


{


Es sei


\mathl{(M,d)}{} ein
\definitionsverweis {metrischer Raum}{}{} und


\mavergleichskette


{\vergleichskette


{ T
}


{ \subseteq }{ M
}


{  }{
}


{    }{
}


{   }{
}


}
{}{}{}
eine nichtleere Teilmenge. Zeige, dass durch


\mavergleichskettedisp


{\vergleichskette


{ d_T(x)
}


{ \defeq} { \inf \left(  d(x,y), \, y \in T  \right)
}


{ } {
}


{ } {
}


{ } {
}


}
{}{}{}
eine wohldefinierte,
\definitionsverweis {stetige Funktion}{}{}
\maabb {} { M } {\R
} {}
gegeben ist.


}


{} {}


\inputaufgabe


{}


{


Es sei


\mathl{(M,d)}{} ein
\definitionsverweis {metrischer Raum}{}{} und


\mavergleichskette


{\vergleichskette


{ T
}


{ \subseteq }{ M
}


{  }{
}


{    }{
}


{   }{
}


}
{}{}{}
eine  Teilmenge. Zeige, dass $T$ genau dann
\definitionsverweis {abgeschlossen}{}{}
ist, wenn es eine
\definitionsverweis {stetige}{}{}
Funktion
\maabb {f} { M } {\R
} {}
mit


\mavergleichskettedisp


{\vergleichskette


{f^{-1}(0)
}


{ =} { T
}


{ } {
}


{ } {
}


{ } {
}


}
{}{}{}
gibt.


}


{} {}


\inputaufgabe


{}


{


Zeige, dass die offenen und die abgeschlossenen Bälle
\mathkor {} {U \left(   P ,r  \right)} {bzw.} {B  \left(  P ,r \right)} {}
im $\R^n$ zu


\mavergleichskette


{\vergleichskette


{r
}


{ > }{ 0
}


{  }{
}


{    }{
}


{   }{
}


}
{}{}{}
nicht offen bzw. abgeschlossen in der
\definitionsverweis {Zariski-Topologie}{}{}
sind.


}


{} {}


\inputaufgabe


{}


{


Charakterisiere in $\Z$ die
Radikale
 mithilfe der Primfaktorzerlegung.


}


{} {}


Ein Element $a$ eines
\definitionsverweis {kommutativen Ringes}{}{}
$R$ heißt \definitionswort {nilpotent}{,} wenn


\mavergleichskette


{\vergleichskette


{ a^n
}


{ = }{ 0
}


{  }{
}


{    }{
}


{   }{
}


}
{}{}{}
für eine natürliche Zahl $n$ ist.


\inputaufgabe


{}


{


Es sei $R$ ein
\definitionsverweis {kommutativer Ring}{}{} und es seien


\mavergleichskette


{\vergleichskette


{ f,g
}


{ \in }{ R
}


{  }{
}


{    }{
}


{   }{
}


}
{}{}{}
\definitionsverweis {nilpotente Elemente}{}{.}
Zeige, dass dann die Summe $f+g$ ebenfalls nilpotent ist.


}


{} {}


\inputaufgabegibtloesung


{}


{


Es sei $R$ ein
\definitionsverweis {kommutativer Ring}{}{} und


\mavergleichskette


{\vergleichskette


{ f
}


{ \in }{ R
}


{  }{
}


{    }{
}


{   }{
}


}
{}{}{}
ein
\definitionsverweis {nilpotentes Element}{}{.}
Zeige, dass $1+f$ eine
\definitionsverweis {Einheit}{}{}
ist.


}


{} {}


\inputaufgabe


{}


{


Es sei $R$ ein
\definitionsverweis {kommutativer Ring}{}{} und $r\in R$ ein
\definitionsverweis {nilpotentes}{}{}
Element. Konstruiere dazu ein lineares Polynom in $R[X]$, das eine
\definitionsverweis {Einheit}{}{}
ist. Man gebe auch das
\definitionsverweis {Inverse}{}{}
dazu an.


}


{} {}


Ein
\definitionsverweis {kommutativer Ring}{}{}
$R$ heißt \definitionswort {reduziert}{,} wenn $0$ das einzige
\definitionsverweis {nilpotente Element}{}{}
von $R$ ist.


\inputaufgabegibtloesung


{}


{


Zeige, dass ein
\definitionsverweis {Ideal}{}{}
${\mathfrak a}$ in einem
\definitionsverweis {kommutativen Ring}{}{}
$R$ genau dann ein
\definitionsverweis {Radikal}{}{}
ist, wenn der
\definitionsverweis {Restklassenring}{}{}
$R/ {\mathfrak a}$
\definitionsverweis {reduziert}{}{}
ist.


}


{} {}


\inputaufgabe


{}


{


Es sei


\mavergleichskette


{\vergleichskette


{ {\mathfrak a}
}


{ \subseteq }{ R
}


{  }{
}


{    }{
}


{   }{
}


}
{}{}{}
ein
\definitionsverweis {Ideal}{}{}
in einem
\definitionsverweis {kommutativen Ring}{}{}
$R$. Zeige, dass die
\definitionsverweis {Potenzen}{}{}


\mathl{{\mathfrak a}^n,\, n \in \N_+}{,} alle dasselbe
\definitionsverweis {Radikal}{}{}
besitzen.


}


{} {}


\inputaufgabe


{}


{


Zeige, dass ein
Primideal
 ein
Radikal
 ist.


}


{} {}


\inputaufgabe


{}


{


Es seien $R$ und $S$
\definitionsverweis {kommutative Ringe}{}{}
und sei
\maabb {\varphi} { R } { S
} {}
ein
\definitionsverweis {Ringhomomorphismus}{}{.}
Es sei ${\mathfrak a}$ ein
\definitionsverweis {Radikal}{}{}
in $S$. Zeige, dass das Urbild


\mathl{\varphi^{-1}( {\mathfrak a} )}{} ein Radikal in $R$ ist.


}


{} {}


\inputaufgabe


{}


{


Es seien ${\mathfrak a}$ und ${\mathfrak b}$
\definitionsverweis {Ideale}{}{}
in


\mathl{K[X_1 , \ldots , X_n]}{} derart, dass ihre
\definitionsverweis {Radikale}{}{}
gleich sind. Zeige, dass dann auch ihre
\definitionsverweis {Nullstellenmengen}{}{}
übereinstimmen. Zeige durch ein Beispiel, dass die Umkehrung nicht stimmt.


}


{} {}


\inputaufgabe


{}


{


Es sei $K$ ein unendlicher Körper und sei


\mavergleichskette


{\vergleichskette


{ F
}


{ \in }{ K[X_1 , \ldots , X_n]
}


{  }{
}


{    }{
}


{   }{
}


}
{}{}{}
ein Polynom. Es sei


\mavergleichskette


{\vergleichskette


{ U
}


{ \subseteq }{ { {\mathbb A}_{ K }^{ n  } }
}


{  }{
}


{    }{
}


{   }{
}


}
{}{}{}
eine nicht-leere Zariski-offene Teilmenge. Es sei


\mavergleichskette


{\vergleichskette


{  F|_{U}
}


{ = }{ 0
}


{  }{
}


{    }{
}


{   }{
}


}
{}{}{}
die Nullfunktion. Zeige, dass dann $F$ das Nullpolynom ist.


}


{} {}


  \zwischenueberschrift{Aufgaben zum Abgeben}


\inputaufgabe


{3}


{


Es sei
\maabbdisp {\varphi} { { {\mathbb A}_{  K }^{  n   } }  } { { {\mathbb A}_{  K }^{  m   } }
} {}
eine Abbildung, die durch $m$ Polynome in $n$ Variablen gegeben sei. Zeige, dass $\varphi$ stetig bezüglich der
\definitionsverweis {Zariski-Topologie}{}{}
ist.


}


{} {}


\inputaufgabe


{4}


{


Es sei $K$ ein unendlicher Körper. Zeige, dass jede nichtleere
\definitionsverweis {Zariski-offene}{}{}
Menge


\mavergleichskette


{\vergleichskette


{ U
}


{ \subseteq }{ { {\mathbb A}_{  K }^{  n   } }
}


{  }{
}


{    }{
}


{   }{
}


}
{}{}{}
\definitionsverweis {dicht}{}{}
ist.


}


{} {Tipp: Reduziere auf den Fall


\mavergleichskette


{\vergleichskette


{ n
}


{ = }{ 1
}


{  }{
}


{    }{
}


{   }{
}


}
{}{}{.}
Verwende nicht

Aufgabe 3.18
!}


\inputaufgabe


{5}


{


Bestimme zu den folgenden Teilmengen der affinen Ebene $\mathbb A^2_K$ den
\definitionsverweis {Zariski-Abschluss}{}{.}
\aufzaehlungfuenf{


\mathl{\left\{  (x, \sin (x))   \mid   x \in \R         \right\}}{,}
}{


\mathl{\left\{  (\cos (x), \sin (x))   \mid   x \in \R         \right\}}{,}
}{


\mathl{\left\{  (x,x^3)   \mid   0 \leq x \leq 1 , \,  x \in \R       \right\}}{,}
}{


\mathl{\left\{  (x,x^3)   \mid   0 \leq x \leq 1 , \,  x \in \Q       \right\}}{,}
}{


\mathl{\left\{  (x,x^3)   \mid   x \in \Z/(5)         \right\}}{.}
}


}


{} {}


Die folgende Aufgabe benutzt einige weiterführende topologische Begriffe.


\inputaufgabe


{4}


{


Es sei $K$ ein
\definitionsverweis {Körper}{}{.}
\aufzaehlungvier{Man zeige, dass für


\mavergleichskette


{\vergleichskette


{ K
}


{ = }{ \R
}


{  }{
}


{    }{
}


{   }{
}


}
{}{}{}
bzw.


\mavergleichskette


{\vergleichskette


{ K
}


{ = }{ {\mathbb C}
}


{  }{
}


{    }{
}


{   }{
}


}
{}{}{}
die Standardtopologie
\zusatzklammer {die metrische oder euklidische Topologie} {} {}
feiner ist als die
\definitionsverweis {Zariski-Topologie}{}{}
auf


\mathl{{\mathbb A}^{1}_{K}}{.}
}{Man zeige, dass für ${\mathbb A}^{1}_{K}$ die Zariski-Topologie mit der
\definitionsverweis {kofiniten Topologie}{}{}
übereinstimmt. Gilt dies auch für


\mathl{{ {\mathbb A}_{ K }^{ n  } }}{} mit


\mavergleichskette


{\vergleichskette


{ n
}


{ \geq }{ 2
}


{  }{
}


{    }{
}


{   }{
}


}
{}{}{?}
}{Wann erfüllt die Zariski-Topologie auf


\mathl{{ {\mathbb A}_{ K }^{ n  } }}{} die Eigenschaft $T_1$, wann ist sie
\definitionsverweis {hausdorffsch}{}{?}
}{Wie sieht die Zariski-Topologie auf


\mathl{{ {\mathbb A}_{ K }^{ n  } }}{} aus, wenn $K$ ein endlicher Körper ist?
}


}


{} {}
